%================================================
\section{Introduction}
\subsection{Contexte}
\begin{frame}{Contexte}
	\begin{itemize}
		\item La distance de décollage est un paramètre essentiel pour la sécurité aérienne.
		\item Elle dépend des caractéristiques aérodynamiques et propulsives de l'avion.
		\item Elle varie selon les conditions atmosphériques et l'état de la piste.
		\item Une estimation fiable est indispensable pour l'exploitation des aéronefs.
	\end{itemize}
	
	\vspace{0.4cm}
	
	\begin{block}{Objectif}
		Développer et comparer plusieurs modèles permettant d'estimer la distance de décollage d'un avion.
	\end{block}
\end{frame}

%================================================
\subsection{Facteurs influençant le décollage}

\begin{frame}{Facteurs influençant la distance de décollage}
	
	\begin{columns}
		
		\column{0.5\textwidth}
		
		\begin{itemize}
			\item Masse de l'avion
			\item Température
			\item Altitude-densité
			\item Vent de face / arrière
			\item Etat de la piste
			\item Effet de sol
		\end{itemize}
		
		\column{0.5\textwidth}
		
		
		
	\end{columns}
	
\end{frame}
